\documentclass{easychair}

\usepackage{amsmath,amssymb}
\usepackage{booktabs}
\usepackage{graphicx}
\usepackage{xspace}
\usepackage{multirow}
\usepackage{url}

\newcommand{\Eprover}{E\xspace}
\newcommand{\Vampire}{Vampire\xspace}
\newcommand{\Mizar}{Mizar\xspace}
\newcommand{\MML}{MML\xspace}
\newcommand{\ArtifactURL}{\textbf{[artifact URL to be inserted before submission]}}

\title{Neural Conjecturing for Saturation Theorem Provers}

\author{
Karel Chvalovsk{\'y}\inst{1}
\and
Martin Suda\inst{1}
\and
Josef Urban\inst{1,2}
}

\institute{
Czech Technical University in Prague
\and
AI4REASON and University of Gothenburg
}

\authorrunning{Chvalovsk{\'y}, Suda, Urban}
\titlerunning{Neural Conjecturing for Saturation Theorem Provers}

\begin{document}
\maketitle

\begin{abstract}
We present a system for generating intermediate lemmas, used as cut
formulas, that restructure first-order proof search in saturation-based
theorem provers.  Given a CNF problem, a heterogeneous graph neural
network encodes clauses, literals, terms, variables and symbols; an
autoregressive decoder then samples candidate lemmas by pointing to
symbols in the input and by obeying arity constraints during decoding.
Candidate lemmas are not trusted: they are accepted only when an ATP
proves both sides of the cut, i.e.\ the problem extended with the
lemma and with its negation.  This gives a sound evaluation loop and a
supervised learning signal.  From the \Mizar Mathematical
Library~\cite{Grabowski2010} we build datasets with both \Eprover and
\Vampire; the \Vampire dataset contains 196,199 useful cut examples
from 2.47M evaluated proof-lemma pairs.  On the \Eprover test set,
the best single configuration helps 96 of 187 baseline-proved
problems, while an ensemble helps 127 (67.9\%), with speedups up to
$196\times$.  More importantly, an \Eprover-trained generator transfers
to \Vampire: with speculative AVATAR claims it proves 582 previously
unsolved \Mizar problems, including 442 that had not been solved by any
ATP in our comparison, and speeds up 2,236 further problems.
\end{abstract}

\section{Introduction}

Conjecturing---the creation of useful intermediate statements---is a
core mathematical ability that complements proof search.  Humans
routinely decompose hard problems into lemmas, and the ability to
propose good lemmas is often what separates a feasible proof from an
intractable search.  Automated conjecturing has been studied since
Lenat's AM~\cite{Lenat1977}, with systems such as Fajtlowicz's
Graffiti~\cite{Fajtlowicz1988} and Colton's HR~\cite{Colton2002}
exploring theory formation via heuristics or controlled enumeration.
More recently, Urban and Jakub\r{u}v~\cite{Urban2020} proposed using
neural sequence models to generate conjectures over large formal
libraries.

In the context of saturation-based theorem proving, conjecturing takes
a specific form: the classical cut rule.  A formula~$L$ splits a hard
problem~$P$ into two subproblems, $P\cup\{L\}$ and $P\cup\{\neg L\}$.
If both are unsatisfiable, then $P$ is unsatisfiable.  Modern provers
such as \Eprover~\cite{Schulz2002} and \Vampire~\cite{Kovacs2013}
combine many sophisticated search heuristics, yet a single unfortunate
trajectory may spend millions of inferences before finding a short
refutation.  A well-chosen cut can bypass such bottlenecks entirely.
The formula~$L$ need not be a user-interesting theorem; it only has to
expose a useful case split.

This paper presents a system for automated cut discovery.  We train a
neural generator on proof lemmas that have been independently measured
to shorten proof search, then use the ATP itself as a checker and
scheduler for generated candidates.  The approach is attractive
for theorem proving because the learned part is entirely speculative:
a bad conjecture may waste time, but cannot make the prover unsound.
The most useful generated conjectures can also be stored as proof-search
hints or as new training examples.

The pipeline is:
\begin{enumerate}\itemsep0pt
\item solve library problems and extract proof lemmas;
\item evaluate each lemma as a cut and keep only speedups;
\item train a graph-to-clause neural generator on the useful cuts;
\item sample diverse conjectures for new problems; and
\item verify them by ATP calls, or in \Vampire by AVATAR claim formulas.
\end{enumerate}
The code, trained checkpoints,
generated conjectures, ATP command lines, logs, and scripts to recreate
the tables are available at: \ArtifactURL.

Our main contributions are as follows.  First, we give a
symbol-independent graph-to-clause architecture for first-order cut
formula generation.  Second, we construct much larger cut datasets than
previous work, including 196K useful \Vampire examples.  Third, we give
a prover-level evaluation rather than a purely syntactic or likelihood
metric.  Finally, we show strong cross-prover transfer: a model trained
from \Eprover proof lemmas helps \Vampire solve hundreds of additional
\Mizar problems.

\section{Learning Cuts from Proofs}
\label{sec:cuts}

The framework for measuring lemma usefulness as proof shortening was
introduced by Goertzel and Urban~\cite{Urban2019}, who created the
first dataset of 3,161 \Mizar CNF problems with cut ratios and
explored GNN-based lemma classification.\footnote{The original
  dataset and code: \url{https://github.com/JUrban/E_conj}}
For a problem $P$ and candidate lemma $L$, let $C_0$ be the baseline
cost of proving $P$, and let $C_1,C_2$ be the costs of proving
$P\cup\{L\}$ and $P\cup\{\neg L\}$.  We define the cut ratio
\[
        r(L,P) = \frac{C_1+C_2}{C_0}.
\]
A pair is useful when both subproofs succeed and $r<1$.  The cost is
processed clauses for \Eprover and burned instructions for \Vampire.
This is the key supervision signal: useful cuts are not labeled by
humans, but by measured search reduction in a sound proof experiment.

\paragraph{AVATAR claims.}
Running two separate proofs is expensive.  \Vampire's AVATAR
architecture~\cite{Voronkov2014} can instead introduce $L$ through the TPTP claim role~\cite{Einarsdottir2024}.
Its SAT layer creates complementary assumptions corresponding to $L$
and $\neg L$, interleaves the two branches, and shares clauses derived
without the claim.  A proof is complete only when the use of $L$ and
the proof of $L$ have both been discharged.  This lets us test generated
conjectures speculatively in a single \Vampire run, and also permits
several generated claims to compete in the same search.

\paragraph{Datasets.}
The \Eprover development data uses 3,161 post-proof-minimized \Mizar
CNF problems exported in TPTP style~\cite{Sutcliffe2009}: \Eprover solves 1,888 and produces 122,356 evaluated
(problem, lemma) pairs.  The larger \Vampire data uses the Deepire~2.0 system~\cite{Suda2025},
which trains GNN-based clause selection inside \Vampire by an
iterative train/eval loop.  A non-AVATAR strategy was first obtained
by hill climbing on a subset of the Mizar40 corpus (3,000 problems,
3,000M instruction limit) and then used with Deepire to train clause
guidance on all 57,880 Mizar40 CNF problems (16,000M limit, 18
iterations, solving 40.5\%$\to$57.7\%).  We use this model to extract
proof lemmas.  For the AVATAR-based evaluation of generated
conjectures, a separate Deepire model was trained with AVATAR enabled,
first on a random half of Mizar40 (20 iterations,
42.1\%$\to$62.0\%), and then on the premise-selected Mizar60 corpus
(22 iterations, 60.1\%$\to$70.6\% on 27,916 problems).  The
non-AVATAR model solves 34,182 of 57,880 full-axiom problems, from
which 2.47M proof-lemma pairs are evaluated; 196,199 are useful cuts
and 3,591 problems have at least one useful lemma.  We also evaluate on
the premise-selected Mizar60 corpus of 55,832 problems, whose median CNF
size is 34 clauses after clausification, giving a more realistic middle
ground between post-proof minimized and full-axiom problems.
Table~\ref{tab:data} summarizes the main corpora.  The useful
cuts are highly skewed in quality: on the \Vampire corpus, 66,461
examples give at least a $2\times$ speedup and 3,190 give at least a
$10\times$ speedup.  We keep all useful examples during training, but
the loss weighting makes the high-speedup tail more influential.

\begin{table}[t]
\centering
\caption{Data and evaluation corpora.  Useful means both cut branches
prove and $r<1$.}
\label{tab:data}
\small
\begin{tabular}{lrrrr}
\toprule
Corpus & Problems & Solved & Evaluated pairs & Useful cuts \\
\midrule
\Eprover minimized CNF & 3,161 & 1,888 & 122,356 & training/eval set \\
\Vampire full-axiom CNF & 57,880 & 34,182 & 2,470,292 & 196,199 \\
Mizar60 premise-selected CNF & 55,832 & 34,080 & ongoing & -- \\
\bottomrule
\end{tabular}
\end{table}

\begin{table}[t]
\centering
\caption{Quality of the 196,199 useful \Vampire cut examples.}
\label{tab:ratios}
\small
\begin{tabular}{rrrrr}
\toprule
$r \leq$ & Examples & Problems & \% examples & \% problems \\
\midrule
0.10 & 3,190 & 164 & 1.6 & 4.6 \\
0.20 & 15,165 & 553 & 7.7 & 15.4 \\
0.50 & 66,461 & 1,645 & 33.9 & 45.8 \\
0.80 & 135,380 & 2,783 & 69.0 & 77.5 \\
0.95 & 179,059 & 3,371 & 91.3 & 93.9 \\
1.00 & 196,199 & 3,591 & 100.0 & 100.0 \\
\bottomrule
\end{tabular}
\end{table}

\section{The Generator}
\label{sec:generator}

\paragraph{Problem graph.}
The input CNF is represented as a heterogeneous graph with five node
types: clause, literal, term, variable, and symbol.  Sixteen directed
edge types encode containment and syntax, including clause-to-literal,
literal-to-argument, term-to-subterm and reverse edges.  Four message
passing layers are applied.  Positional edges use edge-aware GINE
convolutions~\cite{Hu2020}; the remaining edges use GraphSAGE-style
aggregation~\cite{Hamilton2017}.  The final symbol-node embeddings are
both the memory over which the decoder attends and the objects to which
it may point.

\paragraph{Why point to symbols?}
The generator does not emit symbol names from a fixed vocabulary.
Predicates and functions are selected by a pointer distribution over
symbol nodes in the current problem.  Thus the model can be applied to
problems containing symbols not seen in training: it relies on arity,
role, and occurrence structure rather than on names alone.  Optional
named \Mizar embeddings are added for recurring symbols, but unknown
symbols still remain generatable through their graph embeddings.

\paragraph{Decoder and syntax constraints.}
A Transformer decoder generates a clause as actions:
\textsc{NewLitPos}, \textsc{NewLitNeg}, \textsc{Pred},
\textsc{ArgVar}, \textsc{ArgFunc}, \textsc{EndArgs}, and
\textsc{EndClause}.  At each step it predicts an action, a symbol
pointer when the action needs a predicate or function, and a variable
slot when the action needs a variable.  Role masks prevent predicates
from being used as functions and conversely.  An arity stack records
how many arguments remain open; actions inconsistent with the stack are
masked to $-\infty$ before sampling.  This simple state machine is
responsible for the high syntactic validity in the experiments.

\paragraph{Training and variants.}
The loss is the sum of cross-entropies for action, symbol pointer and
variable prediction, ignoring irrelevant outputs at each step.  Samples
are weighted by $1/(1+r)$ so that stronger speedups matter more, and
per-problem normalization prevents a few lemma-rich theorems from
dominating training.  We train Transformer, conditional-VAE and
state-space decoder variants inspired by selective state-space models~\cite{Gu2023}, varying dimension, named-symbol
embeddings and sampling temperature.  Generation uses nucleus sampling
with role masks and removes duplicate clauses before ATP evaluation.

\paragraph{Engineering.}
The implementation has two performance bottlenecks: graph construction
and batching.  For the \Vampire data, a single problem may produce
thousands of graph nodes, and the heterogeneous GNN applies a separate
convolution per edge type.  We therefore precompute graph samples,
cache them in a compact tensor format, and batch by a total-node budget
rather than by a fixed sample count.  A batch of small problems may
contain more than 100 samples, while a batch of large problems may
contain only a few; the loss is normalized by token count, so the
objective is invariant to this packing.  Persistent data-loader workers
and mixed precision keep the A100 GPUs saturated enough for experiments
at the 100K-sample scale.

\section{Experiments}
\label{sec:experiments}

All reported successes are prover-verified.  The generator proposes
formulas, but a result is counted only when the corresponding ATP proof
succeeds.  For \Eprover, we test each generated $L$ by two independent
2-second runs for $P\cup\{L\}$ and $P\cup\{\neg L\}$.  For \Vampire,
we use AVATAR claims with the Deepire~2.0 AVATAR-trained clause
guidance model and a 100,000M instruction limit plus a 30-second
wall-clock limit.

\paragraph{Protocol details.}
The baseline for a problem is always measured without generated
conjectures, using the same prover family and resource accounting as
in the corresponding experiment.  A generated conjecture is discarded
before proving if it is syntactically invalid, uses symbols with
incorrect roles, exceeds the arity-state limits, or duplicates an
earlier sample modulo printing.  For solved baseline problems we report
speedup ratios; for unsolved baseline problems we report newly proved
problems.  This separates two benefits that are often conflated: cuts
can either make an existing proof cheaper, or move a problem from
unsolved to solved within a fixed budget.  All aggregate counts are by
problem, not by lemma, so a problem with many useful conjectures is
counted once.

\subsection{E Prover Results}

The \Eprover test split contains 279 problems, of which 187 are solved
by the 10-second baseline.  We train 15 configurations and generate
40--80 conjectures per problem using nucleus sampling.  Table~\ref{tab:e}
shows selected results.  The best single configuration helps 96 of the
187 baseline-solved problems (51.3\%).  The union of all configurations
helps 127 problems (67.9\%), showing that diversity across models and
sampling regimes is essential.  The best individual speedup is
$196\times$ on \texttt{l101\_xreal\_1}; other large speedups include
$50\times$ on \texttt{t87\_qc\_lang2} and $38\times$ on
\texttt{t35\_xboole\_1}.

\begin{table}[t]
\centering
\caption{Selected \Eprover results on 187 baseline-proved test
problems.  AvgR is the best observed cut ratio averaged over helped
problems; lower is better.}
\label{tab:e}
\small
\begin{tabular}{lrrrr}
\toprule
Configuration & Decoder & Helped & \% & AvgR \\
\midrule
\texttt{a256\_wide}       & Transformer & 96 & 51.3 & 0.558 \\
\texttt{af6\_a\_big\_diverse} & Transformer & 91 & 48.7 & 0.574 \\
\texttt{a256\_bal\_diverse} & Transformer & 87 & 46.5 & 0.524 \\
\texttt{a512\_diverse}    & Transformer & 80 & 42.8 & 0.536 \\
\texttt{af6\_c\_named}    & CVAE        & 78 & 41.7 & 0.597 \\
\texttt{af6\_d\_named}    & SSM         & 72 & 38.5 & 0.612 \\
\texttt{a1024\_diverse}   & Transformer & 67 & 35.8 & 0.530 \\
\midrule
Union of all 15 & mixed & 127 & 67.9 & 0.530 \\
\bottomrule
\end{tabular}
\end{table}

The useful conjectures usually occur early.  Across representative
models, the median rank of the first useful conjecture is between 2
and 5; evaluating only the first five conjectures already recovers a
large fraction of helped problems.  Wider sampling has worse average
rank but better total coverage, which suggests a practical two-phase
strategy: run a cheap top-$k$ filter across several diverse models, and
spend the remaining ATP budget on configurations whose first successes
are complementary.  Table~\ref{tab:ensemble} shows the effect: the best
single model covers 96 problems, but five carefully chosen methods
cover 123, very close to the 127-problem oracle over all 15 methods.

\begin{table}[t]
\centering
\caption{Complementarity on the \Eprover test set.  ``Saved'' is total
baseline processed clauses saved by the best cut for each helped
problem.}
\label{tab:ensemble}
\small
\begin{tabular}{clrr}
\toprule
$N$ & Best combination by coverage & Helped & Saved \\
\midrule
1 & \texttt{a256\_wide} & 96 & 197K \\
2 & + \texttt{d\_named} & 106 & 218K \\
3 & + \texttt{a\_named\_big} & 113 & 286K \\
4 & + \texttt{a256\_bal\_diverse} & 119 & 287K \\
5 & + \texttt{af6\_a\_named} & 123 & 292K \\
\midrule
15 & oracle union & 127 & 301K \\
\bottomrule
\end{tabular}
\end{table}

\begin{table}[t]
\centering
\caption{Effect of evaluating only the top-$k$ conjectures for three
representative \Eprover models.}
\label{tab:rank}
\small
\begin{tabular}{r|rrr|rrr}
\toprule
 & \multicolumn{3}{c|}{Problems helped}
 & \multicolumn{3}{c}{Avg. best ratio} \\
$k$ & \texttt{a256\_wide} & \texttt{a1024} & \texttt{a\_big}
    & \texttt{a256\_wide} & \texttt{a1024} & \texttt{a\_big} \\
\midrule
1  & 20 & 31 & 35 & 0.378 & 0.433 & 0.490 \\
5  & 54 & 53 & 63 & 0.439 & 0.478 & 0.536 \\
10 & 70 & 60 & 75 & 0.503 & 0.498 & 0.551 \\
20 & 81 & 65 & -- & 0.518 & 0.529 & -- \\
50 & 90 & 67 & -- & 0.551 & 0.530 & -- \\
\bottomrule
\end{tabular}
\end{table}

\subsection{Cross-Prover Transfer to Vampire}

The central experiment asks whether cut patterns learned from one
prover are useful to another.  We take an \Eprover-trained $d=256$
model with named symbol embeddings and generate 42 conjectures for each
of 11,881 hard or previously unsolved Mizar60 problems, producing
517,435 conjectures with 98.9\% syntactic validity.  These conjectures
are given to \Vampire as AVATAR claims, so each can be used only if its
proof obligation is also closed.

\begin{table}[t]
\centering
\caption{Cross-prover transfer: \Eprover-trained conjectures evaluated
with \Vampire on hard and unsolved Mizar60 problems.}
\label{tab:vampire}
\small
\begin{tabular}{lr}
\toprule
Problems evaluated & 11,881 \\
Generated conjectures & 517,435 \\
Syntactically valid conjectures & 98.9\% \\
\midrule
Baseline proved by \Vampire, no conjectures & 3,738 \\
\quad of which hard (never solved by any ATP) & 620 \\
Previously unproved, solved with conjectures & 582 \\
\quad of which hard (never solved by any ATP) & 442 \\
Total proved (baseline + conjectures) & 4,320 \\
\quad total hard problems proved & 1,062 \\
Problems with at least 10\% speedup & 2,236 \\
Mean ratio on already-solved helped problems & 0.610 \\
Median rank of first useful conjecture & 3 \\
\bottomrule
\end{tabular}
\end{table}

Table~\ref{tab:vampire} gives the results.  Without generated
conjectures, the AVATAR-guided \Vampire baseline proves 3,738 of the
11,881 problems (including 620 from the hard ``minsub'' set that no
prior ATP had solved).  The generated conjectures prove 582 additional
problems, a 15.6\% relative increase.  The effect is concentrated on
the hardest problems: 442 of the 582 newly proved are from the minsub
set, nearly doubling the number of hard problems solved (620 $\to$
1,062).  Among already solved problems, 2,236 obtain at least a 10\%
search reduction (mean ratio 0.610).

This transfer is nontrivial.  The generator was trained on \Eprover
proof lemmas from post-proof minimized problems, but it helps \Vampire
on premise-selected problems with different clauses, guidance and proof
search.  The most plausible explanation is that the model has learned
structural templates for useful case splits rather than prover-specific
syntactic artifacts.  The pointer architecture is important here: a
fixed vocabulary generator would be much less robust to new symbols and
new premise selections.

\paragraph{Closing the loop.}
The successful \Vampire claims yield 24,025 additional useful examples.
We use them to test whether conjecture generation can become an
iterative data-production loop.  Training from scratch on these examples
overfits quickly.  Initializing from the \Eprover checkpoint is much
better, and the best validation loss is obtained by mixing the original
\Eprover cuts with the new \Vampire cuts (59K combined examples).
The resulting $d{=}256$ model generates 235,089 second-iteration
conjectures (16 per problem) for the same 11,881 problems with
99.3\% syntactic validity.

\Vampire evaluation of the second iteration (Table~\ref{tab:iter2})
shows continued improvement: the cumulative number of hard
(previously ATP-unsolved) problems proved rises from 1,062 (first
iteration, 42 conjectures) to \textbf{1,109} (adding the second
iteration's 16 conjectures)---a 79\% increase over the baseline of
620.  With only 16 conjectures from a model trained on its own
verified outputs, the second iteration adds 47 hard problems beyond
the first, confirming that the loop produces genuinely new and
useful training signal.

\begin{table}[t]
\centering
\caption{Iterative improvement on hard (minsub) \Mizar problems.}
\label{tab:iter2}
\small
\begin{tabular}{lrrr}
\toprule
 & Hard proved & Newly hard & Improvement \\
\midrule
Baseline (no conjectures) & 620 & -- & -- \\
+ 1st iter.\ (42 conj., E-only model) & 1,062 & +442 & +71\% \\
+ 2nd iter.\ (16 conj., combined model) & 1,109 & +489 & +79\% \\
\bottomrule
\end{tabular}
\end{table}

\section{Discussion and Threats to Validity}
\label{sec:discussion}

The experiments support three practical lessons.  First, neural cut
generation should be evaluated by downstream proof search, not by
language-model likelihood.  Many well-formed lemmas are irrelevant, and
some low-probability samples create dramatic speedups.  Second,
generation diversity is a feature rather than noise: the strongest
\Eprover result is an ensemble, and useful \Vampire claims appear
throughout the sampled list.  Third, speculative proving is a natural
interface between learning and ATPs: \Vampire claims let the prover
benefit from a conjecture while simultaneously checking it.

The main threats are cost and selection bias.  Cut evaluation is
expensive, so the datasets inherit the behavior of the prover
configurations used to mine lemmas.  The \Eprover data is small and
post-proof minimized, while the full \Vampire data is larger but harder
to train on because of graph size.  We mitigate these issues with
per-problem weighting, size-aware batching and cross-prover evaluation,
but final generality should be tested on more libraries and on fresh
ATP strategies.

A second threat is budget accounting.  A generated lemma can be
impressively fast once it is known, but a deployment must pay for all
failed candidates as well.  Our rank analysis partially addresses this
by showing that useful cuts often appear early and that model ensembles
can be scheduled greedily, but the current system still relies on ATP
checking as the final ranker.  A learned utility model, trained on the
same cut-ratio labels, is the most direct way to reduce wasted calls.

\paragraph{Artifact design.}
The reviewer artifact is organized around executable claims rather than
around the neural code alone.  It contains: (i) the exact train/validation/test
problem splits; (ii) the generated clauses, including invalid and duplicate
samples; (iii) ATP command lines and resource limits; (iv) raw success and
failure logs; and (v) table-building scripts that recompute all counts from
logs.  This matters because the final theorem-proving claims are not
properties of the neural model in isolation.  They are properties of the
closed loop consisting of sampler, syntactic checker, ATP invocation, timeout
policy and result parser.  The artifact therefore exposes each stage and makes
it possible to rerun smaller slices of the evaluation on a workstation while
also documenting the large cluster runs.

\section{Related Work}

Cut introduction has a long proof-theoretic history beginning with
Gentzen~\cite{Gentzen1935}; automated cut introduction for first-order
logic with equality has been studied by Hetzl et al.~\cite{Hetzl2014}.
Our work differs by generating cuts from a neural model and then
checking them by ATP proof search.  Machine learning for ATPs has most
often focused on premise selection~\cite{KaliszykU2015} and proof or
clause guidance~\cite{Loos2017,Jakubuv2019,Suda2021,Chvalovsky2023}.
Those methods choose what to search next; our method changes the shape
of the problem by adding verified auxiliary formulas.  The closest
antecedents are the cut-usefulness classifier of Goertzel and
Urban~\cite{Urban2019} and the first neural conjecturing datasets of
Urban and Jakub\r{u}v~\cite{Urban2020}; the present system adds a
symbol-pointer graph generator, arity-constrained decoding, large-scale
\Vampire cut data and cross-prover transfer.

\section{Conclusion}

We presented a neural system for generating cut formulas for
saturation-based theorem provers.  The generated formulas are
speculative but soundly verified: a lemma counts only if the ATP
proves both sides of the cut.  On \Eprover, ensembles help 67.9\%
of baseline-proved test problems with speedups up to $196\times$.
On \Vampire, conjectures learned from \Eprover proof lemmas prove
582 additional hard \Mizar problems; a second iteration using
combined training data brings the total of newly proved hard
problems to 489, nearly doubling the baseline of 620 hard problems
solved.  These results suggest that learned cut formulas capture
reusable structure in mathematical proofs and that iterative
generation--evaluation loops can progressively expand the reach of
automated theorem provers.

\begin{thebibliography}{20}

\bibitem{Schulz2002}
S. Schulz.
\newblock E -- a brainiac theorem prover.
\newblock \emph{AI Communications}, 15(2--3):111--126, 2002.

\bibitem{Kovacs2013}
L. Kov{\'a}cs and A. Voronkov.
\newblock First-order theorem proving and Vampire.
\newblock In \emph{CAV 2013}, LNCS 8044, pp. 1--35, 2013.


\bibitem{Suda2025}
M. Suda.
\newblock Efficient Neural Clause-Selection Reinforcement.
\newblock \emph{arXiv:2503.07792}, 2025.

\bibitem{Einarsdottir2024}
S.~H. Einarsd{\'o}ttir, M. Hajd{\'u}, M. Johansson, N. Smallbone, and M. Suda.
\newblock Lemma discovery and strategies for automated induction.
\newblock In \emph{IJCAR 2024}, LNCS 14739, pp. 214--232, 2024.

\bibitem{Voronkov2014}
A. Voronkov.
\newblock AVATAR: The architecture for first-order theorem provers.
\newblock In \emph{CAV 2014}, LNCS 8559, pp. 696--710, 2014.

\bibitem{Gentzen1935}
G. Gentzen.
\newblock Untersuchungen {\"u}ber das logische Schlie{\ss}en.
\newblock \emph{Mathematische Zeitschrift}, 39:176--210, 1935.

\bibitem{Lenat1977}
D.~B. Lenat.
\newblock Automated theory formation in mathematics.
\newblock In \emph{IJCAI 1977}, pp. 833--842, 1977.

\bibitem{Fajtlowicz1988}
S. Fajtlowicz.
\newblock On conjectures of {G}raffiti.
\newblock \emph{Discrete Mathematics}, 72(1--3):113--118, 1988.

\bibitem{Colton2002}
S. Colton.
\newblock Automated Theory Formation in Pure Mathematics.
\newblock Springer, 2002.

\bibitem{Hetzl2014}
S. Hetzl, A. Leitsch, G. Reis, J. Tapolczai, and D. Weller.
\newblock Introducing quantified cuts in logic with equality.
\newblock In \emph{IJCAR 2014}, LNCS 8562, pp. 240--254, 2014.

\bibitem{Urban2019}
Z. Goertzel and J. Urban.
\newblock Usefulness of lemmas via graph neural networks.
\newblock In \emph{AITP 2019}, 2019.

\bibitem{Urban2020}
J. Urban and J. Jakub\r{u}v.
\newblock First neural conjecturing datasets and experiments.
\newblock In \emph{CICM 2020}, LNCS 12236, pp. 315--323, 2020.


\bibitem{Grabowski2010}
A. Grabowski, A. Korni{\l}owicz, and A. Naumowicz.
\newblock Mizar in a nutshell.
\newblock \emph{J. Formalized Reasoning}, 3(2):153--245, 2010.

\bibitem{Sutcliffe2009}
G. Sutcliffe.
\newblock The TPTP problem library and associated infrastructure.
\newblock \emph{J. Automated Reasoning}, 43(4):337--362, 2009.

\bibitem{KaliszykU2015}
C. Kaliszyk and J. Urban.
\newblock MizAR 40 for Mizar 40.
\newblock \emph{J. Automated Reasoning}, 55(3):245--256, 2015.

\bibitem{Loos2017}
S. Loos, G. Irving, C. Szegedy, and C. Kaliszyk.
\newblock Deep network guided proof search.
\newblock In \emph{LPAR-21}, EPiC vol. 46, pp. 85--105, 2017.

\bibitem{Jakubuv2019}
J. Jakub\r{u}v and J. Urban.
\newblock Hammering Mizar by learning clause guidance.
\newblock In \emph{ITP 2019}, LIPIcs vol. 141, pp. 34:1--34:8, 2019.

\bibitem{Suda2021}
M. Suda.
\newblock Vampire with a brain is a good ITP hammer.
\newblock In \emph{FroCoS 2021}, LNCS 12941, pp. 192--209, 2021.

\bibitem{Chvalovsky2023}
K. Chvalovsk{\'y}, K. Korovin, J. Piepenbrock, and J. Urban.
\newblock Guiding an Instantiation Prover with Graph Neural Networks.
\newblock In \emph{LPAR-24}, EPiC vol. 94, pp.~112--123, 2023.

\bibitem{Hu2020}
W. Hu, M. Fey, M. Zitnik, Y. Dong, H. Ren, B. Liu, M. Catasta, and J. Leskovec.
\newblock Open graph benchmark: datasets for machine learning on graphs.
\newblock In \emph{NeurIPS 2020}, 2020.

\bibitem{Hamilton2017}
W. Hamilton, Z. Ying, and J. Leskovec.
\newblock Inductive representation learning on large graphs.
\newblock In \emph{NeurIPS 2017}, 2017.


\bibitem{Gu2023}
A. Gu and T. Dao.
\newblock Mamba: linear-time sequence modeling with selective state spaces.
\newblock \emph{arXiv:2312.00752}, 2023.

\end{thebibliography}

\end{document}
